\documentclass[11pt]{article}
\usepackage{amsmath, amssymb, amsthm}
\usepackage{geometry}
\usepackage{hyperref}
\usepackage{booktabs}
\usepackage{listings}

\geometry{margin=1.1in}

\newtheorem{fact}{Fact}
\newtheorem{rule*}{Rule}

\newcommand{\eps}{\varepsilon}
\newcommand{\dt}{\Delta t}
\newcommand{\dx}{\Delta x}
\newcommand{\code}[1]{\texttt{#1}}
\newcommand{\KE}{\mathrm{KE}}
\newcommand{\R}{\mathbb{R}}

\lstset{
  basicstyle=\ttfamily\small,
  columns=fullflexible,
  keepspaces=true,
  breaklines=true,
  frame=single,
}

\title{The 2D LADMM Solver on a GPU:\\ Compiled Iteration, Parallel Projection, Mixed Precision}
\author{Optimal Transport / Schr\"odinger Bridge Project\\
Companion to the Python/JAX 2D code in \code{python/2d/setup}}
\date{September 2026}

\begin{document}
\maketitle

\begin{abstract}
This note records what was changed in the 2D Python/JAX solver to make it fast on an
NVIDIA GH200, why each change matters, what was measured, and which precision choices
are safe. It is written so the same decisions can be carried to the 3D code. Sections
\ref{sec:model}--\ref{sec:memory} cover speed and memory in double precision; Section
\ref{sec:precision} covers single and mixed precision, including the one non-obvious
result: the algorithm tolerates float32 storage of a variable only if the optimality
condition attached to that variable is evaluated at the stored value. Section
\ref{sec:node} describes the DeltaAI node, which parts of the code run where, what is
serial, and when several GPUs would pay off.
\end{abstract}

\section{The algorithm and the cost model}\label{sec:model}

One LADMM iteration for $\min_{x,y}\, I_{\mathrm{FP}}(x) + \KE(y)$ s.t.\ $Ax = y$ is
\begin{align}
  r^k &= A x^k - y^k, \label{eq:r}\\
  x^{k+1} &= P\!\left(x^k - \tfrac{1}{\tau} A^{\!\top}(\gamma r^k - \delta^k)\right),
  \label{eq:x}\\
  y^{k+1} &= \operatorname{prox}_{\KE/\gamma}\!\left(A x^{k+1} - \delta^k/\gamma\right),
  \label{eq:y}\\
  \delta^{k+1} &= \delta^k - \gamma\,(A x^{k+1} - y^{k+1}), \label{eq:delta}
\end{align}
where $P$ is the projection onto the discrete Fokker--Planck constraint (one DCT in
$(x,y)$, a tridiagonal solve in $t$ per spatial mode, one inverse DCT), $A$ is the
staggered-to-cell-centre interpolation, and the prox of $\KE(\rho,m)=|m|^2/(2\rho)$ is a
pointwise cubic. Every operation is either a transform, a banded solve, or an elementwise
map over arrays of $n_t n_x n_y$ entries; nothing is a dense matrix.

On a GPU the cost of such a loop has three parts, and it is worth knowing which one you
are paying:
\begin{enumerate}
  \item \textbf{Launch latency.} Every kernel costs 5--20 $\mu$s to dispatch, and a
  device-to-host copy (a \code{float()} on an array) stalls the pipeline. An iteration
  written as $\sim$300 small \code{jnp} operations pays $\sim$300 launches whether the
  arrays hold $10^3$ or $10^7$ entries.
  \item \textbf{Memory traffic.} Elementwise maps, interpolations and reductions are
  bound by bytes moved, not flops. A GH200 moves about 4 TB/s; a float64 array of
  $256^3$ entries is 134 MB, so each pass over the state costs $\sim$35 $\mu$s.
  \item \textbf{Arithmetic.} Only the DCTs (as dense matrix products) and the cubic in
  the prox have appreciable flop counts. On H100/GH200 the float64 tensor-core rate
  (67 TFLOP/s) equals the non-tensor float32 rate, so a float64 matmul is not slower
  than a float32 one at the precision this code needs (Section \ref{sec:dct}).
\end{enumerate}
The eager code was entirely in regime 1: 215 ms per iteration at $32^3$ and 185 ms at
$128^3$, i.e.\ flat in the grid size (Table \ref{tab:speed}). Everything in Sections
\ref{sec:jit}--\ref{sec:dct} is about moving from regime 1 to regime 2.

\section{Compiling the iteration}\label{sec:jit}

\code{ladmm\_jit.py} traces one whole iteration \eqref{eq:r}--\eqref{eq:delta}, including
the seven residual and scale norms used by the stopping test, into a single XLA
executable, and runs \code{check\_every} iterations back to back inside a
\code{lax.scan}. The only host synchronisation is the copy of the $k\times 9$ array of
scalar diagnostics at the end of each chunk. Four details make this work.

\paragraph{Everything traceable.} No \code{float()} inside the loop: the norm callbacks
return \code{jnp} scalars, and the stopping test runs on the host on the chunk's
diagnostics. \code{State2D} and the projection's precompute dataclasses are registered as
pytrees (\code{jax.tree\_util.register\_dataclass}) so they pass through \code{jit} and
\code{scan} as arguments rather than being baked in as constants.

\paragraph{One fewer $A$ per iteration.} The residual \eqref{eq:r} equals the $D$ of the
previous iteration, so it is carried instead of recomputed. Bitwise the same numbers.

\paragraph{Diagnostics are not free.} Each norm is a full reduction over the state; on a
GPU the seven of them cost about as much as the KE prox. \code{diag\_level="stopping"}
drops the three descriptive ones ($\|x^{k+1}-x^k\|$ etc.).

\paragraph{Recompilation.} The compiled chunk closes over the pipeline's callables, so
every call to \code{discretize\_then\_optimize\_*} recompiles (0.4--0.6 s on the GH200).
A benchmark that times a fresh 50-iteration solve reports $\sim$10 ms/iteration of pure
compile time. Measure with \code{walltime - walltime\_first\_chunk}, as
\code{bench\_2d.py} does.

With \code{check\_every=1} the iterates equal the eager solver's to rounding (XLA may
fuse $a b + c$ into an FMA). With \code{check\_every=k} the stop is detected up to $k-1$
iterations late; every extra iteration is a further LADMM step, so this only moves the
answer closer.

\section{The projection on a GPU}\label{sec:proj}

\subsection{The time solve: sequential Thomas vs.\ parallel DCT}\label{sec:timesolve}

After the DCT in $(x,y)$ the projection is $n_x n_y$ independent tridiagonal systems in
$t$, one per mode $(k_x,k_y)$, all with the same structure: constant interior diagonal,
constant off-diagonal $e$, and two deviating end rows. The MATLAB and the eager Python
code solve them with a batched Thomas sweep, vectorised over modes. Written as
\code{lax.scan} (\code{thomas\_batch\_precomp/solve}) rather than a Python list, the
sweep is $2 n_t$ tiny sequential kernels per projection; on a GPU this is latency-bound
again and is the reason the uniform-grid projection was 2--3$\times$ slower than it
needed to be.

The parallel alternative (\code{method="dct"}, \code{precomp\_dct\_modes}) uses the fact
that each per-mode matrix is
\[
  T = T_0 + a_0\, e_0 e_0^{\!\top} + a_1\, e_{n_t-1} e_{n_t-1}^{\!\top},
\]
with $T_0$ a Neumann Laplacian in $t$, diagonalised by the DCT-II in $t$ with eigenvalues
$\mu_j = \mathrm{mid} + 2e\cos(\pi j/n_t)$, and $(a_0,a_1)$ the two end-row deviations. By
the Woodbury identity \cite{hager1989},
\[
  T^{-1} = T_0^{-1} - T_0^{-1} U\,(I + C\,V^{\!\top} T_0^{-1} U)^{-1} C\, V^{\!\top}
  T_0^{-1},
\]
with $U=V=[e_0\; e_{n_t-1}]$ and $C=\mathrm{diag}(a_0,a_1)$: a $2\times 2$ capacitance
matrix per mode, precomputed once. A solve is then one DCT in $t$, a division by $\mu_j$,
one inverse DCT, and a rank-2 elementwise correction: no sequential dependency at all.
The DC mode $(0,0)$ needs no special case: there $a_0=a_1=0$ and $\mu_0 = 0$, and
zeroing $1/\mu_0$ is exactly the pseudo-inverse the Thomas path applied by hand. The
non-uniform time grid is not Toeplitz in $t$ and keeps the scan-Thomas path; it is
correspondingly slower (Table \ref{tab:speed}).

\subsection{The DCT kernel}\label{sec:dct}

JAX has no DCT. Three implementations of the orthonormal DCT-II along one axis are
provided (\code{set\_dct\_method}): mirror-and-FFT of length $2N$, Makhoul's $N$-point
reordering \cite{makhoul1980}, and a dense $N\times N$ cosine matrix applied with one
\code{tensordot}. On the GH200 the matrix form is the fastest for $N\le 512$ by
1.3--1.9$\times$: XLA's FFT along a non-last axis inserts transposes, whereas the matmul
is a single cuBLAS call on an axis of the caller's choosing. Its cost is $O(N)$ per
entry instead of $O(\log N)$, which does not matter at these sizes. In float32 the matmul
precision must be pinned (\code{jax\_default\_matmul\_precision="highest"}) or it silently
uses TF32 tensor cores with a 10-bit mantissa.

\section{Measured speed}\label{sec:measure}

\begin{table}[h]
\centering
\begin{tabular}{lrrrr}
\toprule
grid $n_t=n_x=n_y$ & eager & jit, scan-Thomas & jit, DCT modes & clustered (scan) \\
\midrule
$32^3$  & 215 & 0.25 & 0.09 & 0.6 \\
$128^3$ & 185 & 2.5 & 0.87 (CN) / 1.04 (ETD) & 3.4 \\
$256^3$ & --  & 14.8 & 6.0--7.6 & 12.3 \\
\bottomrule
\end{tabular}
\caption{Milliseconds per LADMM iteration, float64, GH200, steady state (compile
excluded), \code{check\_every}=50--100. ``DCT modes'' uses the matmul DCT. The clustered
time grid cannot use the DCT-in-$t$ solver.}
\label{tab:speed}
\end{table}

Ordered by effect: compiling the loop (100--2000$\times$) $\gg$ DCT modes vs.\ scan-Thomas
(2--3$\times$ at $32^3$--$128^3$) $>$ matmul DCT vs.\ FFT (1.3--1.9$\times$ at $\ge 128^3$)
$>$ \code{check\_every} 10--50 vs.\ 1 (2$\times$ at $32^3$, 10\% at $128^3$) $>$
\code{diag\_level="stopping"} (5--10\%). The 28-run grid-convergence suite (Gaussian,
7 values of $\eps$, $n=32..256$, 40k--150k iterations) takes 159 min of solve time in
float64.

\section{Memory footprint}\label{sec:memory}

The solver holds about 67 arrays of size $n_t n_x n_y$ in float64 (9.0 GB at $256^3$,
134 MB per array). They are: the LADMM carry $(x, y, By, \delta, D)$ with three
components each (15), the previous $x$ inside the scan (3), the projection's temporaries
(the $(n_t+1)$-plane DCT of $\rho$, the residual, the mode solution, its DCT in $t$, the
weighted potential and its derivatives: $\sim$10), the precompute ($1/\mu$, $g_0$, $g_1$,
$D_{\mathrm{mod}}$: 4), the reference solutions for error tracking (3 per reference), the
diagnostics' temporaries, and XLA's own scratch. Only the first three groups are
essential; the references can be kept off the device and, with $\alpha=1$, the previous
$y$ and $By$ are never read. This is the accounting behind the 3D estimates in
Section \ref{sec:3d}.

\section{Precision}\label{sec:precision}

Float32 halves every byte in Section \ref{sec:memory}. It is also where the algorithm
turns out to be fragile, in a specific and instructive way.

\subsection{Plain float32}

Same 28 runs, everything float32 (setup, precompute and references built in float64 and
cast): 1.05$\times$ faster at $32^3$, 1.5--1.6$\times$ at $\ge 128^3$ uniform,
1.2--1.3$\times$ clustered (the sequential sweep does not care about bytes), 2.2$\times$
less memory. Accuracy: identical to float64 for $n\le 64$; at $n=256$ and small $\eps$
the density error floors at $10^{-3}$ instead of $3\times10^{-5}$; mass conservation
degrades from $10^{-13}$ to $10^{-4}$--$10^{-2}$; on the clustered grid at $n=256$ the
solve diverges for $\eps\ge 0.1$.

\subsection{What needs float64: three hypotheses, one survivor}

\paragraph{The projection.} $R$ and $R^{\!\top}$ are time differences divided by $\dt$.
The residual's rounding is harmless (the solve smooths it: the $1/j^2$ tail of the time
spectrum sums to a constant), but $R^{\!\top}$ applied to the large, smooth potential
turns float32 rounding into an $O(u_{32}\,|\phi|/\dt)$ error, and on the clustered grid
the $N$-weighted per-mode systems span $1/(\dt_{\min}^2\dx^2)$ to $1/\dt_{\max}^2$.
Running the projection in float64 on a float32 iterate (\code{precision.mixed\_projection}:
cast up, project, cast down) restores mass conservation to $10^{-9}$ and removes the
divergence -- but the error floor stays: $6.1\times10^{-4}$ at $n=128$ against float64's
$3.4\times10^{-4}$, and $4.3\times10^{-4}$ at $n=256$ against $3.2\times10^{-5}$, with
the residuals $D$, $P$ identical to float64's.

\paragraph{Sub-ulp steps.} Late LADMM steps are $P/\tau\sim10^{-8}$ relative, below
float32's unit roundoff $u_{32}=6\times10^{-8}$, so a float32 $x$ cannot represent them.
Keeping $x$ and $\delta$ as float64 ``master copies'' (\code{LadmmConfig.master\_dtype}),
the recipe of mixed-precision training \cite{micikevicius2018}, changed nothing:
$5.6\times10^{-4}$ at $n=128$.

\paragraph{Frozen inconsistency in the KKT system.} Two diagnostic runs, each float64
everywhere except for one injected rounding (\code{--precision fp64rx}, \code{fp64px}):
rounding $x$ to float32 before $A$ gives $3.377\times10^{-4}$, float64's value to four
digits; running only the KE prox and $y$ in float32 gives $5.98\times10^{-4}$, the whole
floor. The reason is the identity, exact for the prox of a differentiable function,
\begin{equation}
  \delta^{k+1} = \delta^k - \gamma\,(A x^{k+1} - y^{k+1}) = -\nabla\KE\!\left(y^{k+1}\right),
  \label{eq:dualid}
\end{equation}
which holds at the \emph{exact} prox output. When $y$ is stored in float32 the dual is
the KE gradient at a point $u_{32}$ away from the $y$ the constraint $Ax=y$ actually
uses. That is a frozen, non-random inconsistency in the KKT system, and it is amplified
by the ill-conditioned directions of the problem -- the same slow modes that make LADMM
need $10^5$ iterations -- by roughly $10^3$. The excess error sits in the core of the
density, mid-trajectory, with the shape of the slowest mode, not in the tails. Rounding
$x$ has no such inconsistency: the prox is evaluated at the same rounded $x$ the
constraint sees, so the KKT conditions hold at the stored iterate and the rounding only
selects a nearby solution.

\begin{rule*}
A variable may be stored in reduced precision only if every optimality condition
involving it is evaluated at the stored value. Storage precision is a property of the
fixed point, not of the arithmetic.
\end{rule*}

\subsection{The working scheme (\code{--precision mixedc})}

\begin{itemize}
  \item float32: $y$, $By$, $D$, the consensus map $A$ applied to the rounded $x$, the
  reference-error diagnostics.
  \item float64: $x$, $\delta$, the projection, $A^{\!\top}$, and the \emph{arithmetic}
  of the KE prox (its output is stored float32). The float32 Cardano root is up to
  $7\times10^{-4}$ off in near-empty cells; in float64 the casts fuse into the
  elementwise kernel and cost nothing.
  \item The dual is computed from the stored $y$ by \eqref{eq:dualid}:
  $\delta^{k+1} = -\nabla\KE(y^{k+1}_{32})$ in float64, with the standard update kept
  in the cells the prox clamps to $\rho=0$ (\code{ladmm\_jit.py}, \code{dual\_fn}).
\end{itemize}

Results over the 28 runs: every error agrees with float64 to four digits, the distance
between the two solutions is at most $6\times10^{-7}$, fitted orders 3.68 / 3.68 / 2.52 /
2.02 / 2.05 / 1.87 / 1.02 for $\eps=0,10^{-8},\dots,1$ (float64's), mass conservation
float64's. Speed is 0.8--1.1$\times$ float64 (1.04$\times$ at $256^3$); peak memory is
60\% of float64 at $128^3$ and 71\% at $256^3$ (6.4 GB vs.\ 9.0 GB). The scheme is a
memory tool, not a speed tool: only the cheap, memory-bound parts became float32, and
the float64 matmuls were never the bottleneck.

\subsection{Half precision}\label{sec:half}

Two directions are common in the literature: compute in a low precision and accumulate
in a higher one (mixed-precision training \cite{micikevicius2018}: float16 arithmetic,
float32 master weights), and factorise in a low precision and refine the residual in a
higher one (iterative refinement in three precisions \cite{carson2018,higham2022}).
Neither helps here.
\begin{itemize}
  \item The compute-low direction is what Sections above tested with float32 as the low
  precision: the fixed point moved by $u\times10^3$. With float16 ($u=5\times10^{-4}$) or
  bfloat16 ($u=4\times10^{-3}$) even the consistent scheme would place the solution
  $\sim u$ away from the true one, i.e.\ $10^{-3}$--$10^{-2}$ relative, above the
  discretization error at every $n\ge 64$. The residual $D = Ax - y$ is a cancellation
  of $O(64)$ numbers; in half precision it is noise at the $10^{-2}$ level and the dual
  update \eqref{eq:delta} integrates that noise. Float16's range ($6\times10^{-5}$ to
  $6.5\times10^4$ normal) also flushes the density's tails to zero.
  \item The refine direction needs a cheap low-precision approximate solve and an exact
  residual. The projection is already a direct $O(n)$ solve; a half-precision version
  plus a float64 residual costs more than the float64 solve it would replace.
  \item Half-precision DCTs on tensor cores would put a $10^{-3}$ relative error into the
  projection's residual; this is why even TF32 is disabled.
  \item The memory argument does not hold either: the float32 arrays in \code{mixedc}
  are $\sim$9 of $\sim$48 array-equivalents. Halving them saves 5\%. The bytes are in
  the float64 state, the projection's temporaries and the diagnostics.
\end{itemize}

\section{The DeltaAI node: where the work runs}\label{sec:node}

\subsection{Hardware}

A \code{ghx4} node is four NVIDIA GH200 superchips. Each superchip is a 72-core Grace
CPU (Arm Neoverse-V2, $\sim$120 GB LPDDR5X) and an H100 GPU with 98 GB of HBM3
(\code{nvidia-smi} reports 97{,}871 MiB usable), joined by NVLink-C2C at 900 GB/s. The
four GPUs of a node are connected by NVLink; nodes are connected by HPE Slingshot,
more than an order of magnitude slower than NVLink. A job asks for a share of one node:
the suite runs used \code{--gpus=1 --cpus-per-task=16 --mem=96g}, i.e.\ one GPU and
16 of the 288 cores. The login nodes are the same Grace CPUs without a GPU.

\subsection{What runs on the CPU and what on the GPU}

\begin{itemize}
  \item \textbf{CPU (Python, NumPy, SciPy):} the driver script, the problem setup, the
  1D Sinkhorn reference fits and their spectral resampling, the stopping test on each
  chunk's diagnostics, the per-time-row error rows and the \code{.npz} output after
  the solve. All of this is a few seconds per run; the solve itself is minutes.
  \item \textbf{GPU (XLA):} the precompute (built with \code{jnp}, so it already lives on
  the device), the initial iterate, and every operation inside the compiled chunk:
  $A$, $A^{\!\top}$, both projections, the KE prox, the dual update, the norms, the
  reference errors. The Python interpreter's only job during a solve is to call the
  chunk \code{max\_iter/check\_every} times and read back a $k\times 9$ array.
\end{itemize}
The same compiled code runs on the 72-core Grace by setting \code{--device cpu}
(Table \ref{tab:cpugpu}); the GPU is 20--60$\times$ faster, and the gap grows with the
grid because the CPU is bound by DRAM bandwidth ($\sim$500 GB/s) against HBM3's 4 TB/s.

\begin{table}[h]
\centering
\begin{tabular}{lrrr}
\toprule
grid & Grace CPU, 72 threads & one GH200 GPU & ratio \\
\midrule
$64^3$ uniform, DCT modes   & 5.7  & 0.27 & 21 \\
$128^3$ uniform, DCT modes  & 36.7 & 1.15 & 32 \\
$256^3$ uniform, DCT modes  & 423  & 7.3  & 58 \\
$64^3$ clustered, Thomas    & 5.9  & 1.25 & 5 \\
$128^3$ clustered, Thomas   & 39.6 & 3.07 & 13 \\
\bottomrule
\end{tabular}
\caption{Milliseconds per iteration, ETD (CN for the clustered rows), float64, same
compiled code on both devices (\code{cpu\_vs\_gpu.sbatch}). The sequential Thomas sweep
hurts the GPU more than the CPU: the ratio drops from 32 to 13 at $128^3$.}
\label{tab:cpugpu}
\end{table}

\subsection{Serial and parallel parts}

Within one iteration everything is parallel over the grid except two things, and the
iteration count itself is serial.
\begin{itemize}
  \item \textbf{Parallel over all $n_t n_x n_y$ entries:} $A$, $A^{\!\top}$, the
  residual and the corrections inside the projection, the KE prox, the dual update.
  Each is one or a few fused elementwise kernels.
  \item \textbf{Parallel over $n_x n_y$ modes, sequential over $t$:} the Thomas sweep
  ($2n_t$ dependent steps; the clustered grid). The DCT-in-$t$ solver removes this:
  it is parallel over all entries.
  \item \textbf{Transforms:} the DCTs are dense matrix products, parallel over
  everything, executed by cuBLAS.
  \item \textbf{Reductions:} the norms are tree reductions over the state; parallel,
  but each one is a full pass over memory and ends in a single scalar.
  \item \textbf{Serial by nature:} the LADMM iterations. Nothing about iteration $k+1$
  can start before iteration $k$ ends, so the wall time is
  $\text{iterations}\times\text{ms/iteration}$ and the only levers are the per-iteration
  cost (this note) and the iteration count ($\gamma$, the time grid, the initial
  guess).
\end{itemize}

\subsection{Communication}

With one GPU there is no device-to-device traffic. Host--device traffic per solve: the
initial upload of the problem and precompute (tens of MB), one download of
$9\cdot\code{check\_every}$ doubles per chunk (a few kB), and the final download of the
solution. On a GH200 that path is NVLink-C2C at 900 GB/s, so even the eager code's
per-iteration syncs were latency, not bandwidth. The suite's other communication is the
file system: 2--6 MB of \code{.npz} per run.

\subsection{Multiple GPUs}

The four GPUs of a node hold 390 GB together. Splitting the arrays across them buys
memory and bandwidth; it costs communication at every operator that couples the
distributed axis.

\paragraph{Decomposition.} Slabs along $t$ are the natural choice: the spatial DCTs are
then entirely local, the time-direction operators ($A_t$, $A_t^{\!\top}$, $R_t$,
$R_t^{\!\top}$) need a one-plane halo exchange (16 MB at $128^3$ in space), and only
the per-mode time solve needs the whole $t$ axis on one device: one all-to-all
transpose of the mode array into $t$-complete pencils, solve, transpose back. Slabs
along $z$ move the same transpose to the $z$-DCT instead. The norms need an
all-reduce of one scalar each. In JAX this is \code{jax.sharding} with
\code{shard\_map}; XLA emits the NCCL collectives.

\paragraph{Cost.} Per iteration, two transposes of one float64 array plus halos. At
$128^4$ an array is 2.1 GB; over NVLink the two transposes take on the order of 10--15
ms against $\sim$40 ms of local compute per GPU, so about $3\times$ on four GPUs and
the full node's memory. At 2D sizes ($256^3$: 134 MB per array, 7 ms per iteration)
the transposes are $\sim$1 ms each and the launch overhead of the collectives is
comparable, so the gain would be $1.5\times$ at best.

\paragraph{Verdict.} Not useful for 2D: every 2D grid in this study fits and runs
in minutes on one GPU. Useful in 3D for one reason, memory: $128^4$ in float64 or
\code{mixedc} does not fit on one GPU (Section \ref{sec:3d}) and fits comfortably on
four. Stay within one node; across nodes the transposes would run over Slingshot and
the iteration would be communication-bound. The engineering cost is moderate (the
sharding annotations and one transpose in the projection) and should come after the
3D code has the single-GPU path of Sections \ref{sec:jit}--\ref{sec:proj}.

\section{What carries to 3D}\label{sec:3d}

The 3D tree (\code{python/3d}) still has the eager, list-based Thomas sweep and no
compiled path. In order: register the 3D state and precompute as pytrees; port
\code{ladmm\_jit.py} unchanged (it is grid-agnostic); write the 3D DCT as three
\code{tensordot} calls and the DCT-in-$t$ mode solver with the same Woodbury correction
(only $\lambda_{xy}\to\lambda_{xyz}$ changes); build all zero planes and initial iterates
in \code{problem.rho0.dtype}; then enable \code{mixedc}. Per array, $n^4\times 8$ bytes;
with $\sim$55 array-equivalents (four state components) the single-GPU limit on a 98 GB
GH200 is $112^4$ in \code{mixedc} and $96^4$ in float64; $128^3$ in space fits with
$n_t=64$ or after dropping the references and the unused carry entries.

\begin{thebibliography}{9}
\bibitem{jaxdocs} JAX documentation: \emph{How to think in JAX}, \emph{JAX -- The Sharp
Bits}, \code{jax.lax.scan}, \code{jax.jit}. \url{https://jax.readthedocs.io}.
\bibitem{hager1989} W.~W. Hager, Updating the inverse of a matrix, \emph{SIAM Review}
31 (1989) 221--239.
\bibitem{makhoul1980} J.~Makhoul, A fast cosine transform in one and two dimensions,
\emph{IEEE Trans.\ Acoust.\ Speech Signal Process.} 28 (1980) 27--34.
\bibitem{williams2009} S.~Williams, A.~Waterman, D.~Patterson, Roofline: an insightful
visual performance model for multicore architectures, \emph{Comm.\ ACM} 52 (2009) 65--76.
\bibitem{zhang2010} Y.~Zhang, J.~Cohen, J.~D. Owens, Fast tridiagonal solvers on the
GPU, \emph{PPoPP} 2010.
\bibitem{goldberg1991} D.~Goldberg, What every computer scientist should know about
floating-point arithmetic, \emph{ACM Computing Surveys} 23 (1991) 5--48.
\bibitem{higham2002} N.~J. Higham, \emph{Accuracy and Stability of Numerical Algorithms},
2nd ed., SIAM, 2002.
\bibitem{micikevicius2018} P.~Micikevicius et al., Mixed precision training,
\emph{ICLR} 2018, arXiv:1710.03740.
\bibitem{carson2018} E.~Carson, N.~J. Higham, Accelerating the solution of linear
systems by iterative refinement in three precisions, \emph{SIAM J.\ Sci.\ Comput.} 40
(2018) A817--A847.
\bibitem{higham2022} N.~J. Higham, T.~Mary, Mixed precision algorithms in numerical
linear algebra, \emph{Acta Numerica} 31 (2022) 347--414.
\bibitem{abdelfattah2021} A.~Abdelfattah et al., A survey of numerical linear algebra
methods utilizing mixed-precision arithmetic, \emph{Int.\ J.\ High Perf.\ Comput.\ Appl.}
35 (2021) 344--369.
\end{thebibliography}

\end{document}
